@@M0@@
@@M1@@
defining on the functor @@M2@@ a multiplication and operations of @@M3@@. By 1.13, the multiplication is defined by its value on the element @@M4@@ of @@M5@@; it will be denoted by @@M6@@; we therefore have:
@@M7@@
The @@M8@@-th universal polynomial of this law is the isobaric component of weight @@M9@@ of @@M10@@ (1.15), and will be denoted by @@M11@@.

The operations of @@M12@@ will again be denoted by exponentiation (3.8). The universal polynomials defining these operations are the isobaric components of lowest degree of the @@M13@@ (1.15); since these are isobaric, they are the @@M14@@ themselves.

Let @@M15@@; the ring @@M16@@ defined by the preceding laws is in general without a unit; we therefore extend this ring by adjoining the scalars of @@M17@@, introducing the new ring @@M18@@, endowed with the composition laws defined by: @@M19@@, @@M20@@,
@@M21@@
@@M22@@

The ring thus constructed admits (1,1) as its unit, and @@M23@@ is identified with the ideal @@M24@@. It is naturally endowed with a @@M25@@-algebra structure such that:
@@M26@@

